% Einheit 1 – Wurzelziehen und Lösbarkeit – Nr. 12–23
\setcounter{aufgabe}{11}
\einheitenkopf[e1]{Einheit 1 von 4 · Wurzelziehen und Lösbarkeit}
\zweigzeile{Hier lernst du, Gleichungen mit $x^2$, aber ohne $x$-Glied durch Wurzelziehen zu lösen und zu sagen, wie viele Lösungen sie haben · neu oder schon bekannt – je nach Buch (an der Oberschule je nach Buch Kl. 9 oder 10) · P10 · baut auf: Quadrieren, Wurzelziehen, lineare Gleichungen, Scheitelpunktform}

\begin{aufgabe}{Ich erkenne, ob eine Gleichung quadratisch ist.}
Rechne nicht. Kreuze an, ob die Gleichung quadratisch oder linear ist.
\begin{teile}
\teil $3x + 4 = 10$ \hfill \kreuz{quadratisch} \quad \kreuz{linear}
\teil $x^2 = 50$ \hfill \kreuz{quadratisch} \quad \kreuz{linear}
\teil $(x - 1)^2 = 9$ \hfill \kreuz{quadratisch} \quad \kreuz{linear}
\teil $2\cdot(x + 3) = 8$ \hfill \kreuz{quadratisch} \quad \kreuz{linear}
\teil $5x = x^2$ \hfill \kreuz{quadratisch} \quad \kreuz{linear}
\end{teile}
\end{aufgabe}

\begin{aufgabe}{Ich erkenne an der Zahl rechts, wie viele Lösungen eine Gleichung wie $x^2 = 5$ hat.}
Rechne nicht. Kreuze an, wie viele Lösungen die Gleichung hat.
\begin{teile}
\teil $x^2 = 16$ \hfill \kreuz{zwei} \quad \kreuz{eine} \quad \kreuz{keine}
\teil $x^2 = -9$ \hfill \kreuz{zwei} \quad \kreuz{eine} \quad \kreuz{keine}
\teil $x^2 = 2$ \hfill \kreuz{zwei} \quad \kreuz{eine} \quad \kreuz{keine}
\teil $5x^2 = 0$ \hfill \kreuz{zwei} \quad \kreuz{eine} \quad \kreuz{keine}
\teil $x^2 = -0{,}5$ \hfill \kreuz{zwei} \quad \kreuz{eine} \quad \kreuz{keine}
\end{teile}
\end{aufgabe}

\verfahren{Gleichungen der Form $x^2 = $ Zahl}

\begin{aufgabe}{Ich kann eine Gleichung der Form $x^2 = $ Zahl durch Wurzelziehen lösen.}
Löse die Gleichung.
\rechenplatz{2}
\begin{gleichungsraster}[2]
\gl{x^2 = 25} & \gl{x^2 = 81} \\
\gl{x^2 = 4} & \gl{x^2 = 144} \\
\anweisung{Rechne mit dem Taschenrechner, wenn die Wurzel nicht aufgeht. Gib die Lösungen mit Wurzel und auf zwei Stellen nach dem Komma gerundet an.}
\gl{x^2 = 11} & \gl{x^2 = -16} \\
\end{gleichungsraster}
\end{aufgabe}

\begin{aufgabe}{Ich kann erst $x^2$ freistellen und dann die Wurzel ziehen.}
Löse die Gleichung.
\rechenplatz{3}
\begin{gleichungsraster}[3]
\gl{x^2 + 3 = 19} & \gl{x^2 - 20 = 5} \\
\gl{5x^2 = 45} & \gl{3x^2 + 5 = 32} \\
\gl{4x^2 - 9 = -9} & \gl{2x^2 + 11 = 3} \\
\gl{3x^2 - 4 = 17} & \\
\end{gleichungsraster}
\end{aufgabe}

\verfahren{Gleichungen mit quadrierter Klammer}

\begin{aufgabe}{Ich kann eine Gleichung mit quadrierter Klammer durch Rückwärtsrechnen lösen.}
Löse die Gleichung.
\rechenplatz{3}
\begin{gleichungsraster}[3]
\gl{(x - 2)^2 = 9} & \gl{(x - 1)^2 = 16} \\
\gl{(x - 5)^2 = 4} & \gl{(x - 3)^2 = 25} \\
\gl{(x + 3)^2 = 49} & \gl{(x + 6)^2 = 0} \\
\gl{(x - 2)^2 = 5} & \\
\end{gleichungsraster}
\end{aufgabe}

\verfahren{Lösungen prüfen und zählen}

\begin{aufgabe}{Ich kann prüfen, ob eine Zahl eine Lösung einer quadratischen Gleichung ist.}
Setze die Zahl ein und kreuze an.
\begin{teile}
\teil $x = 4$ in $x^2 - 16 = 0$ \hfill \kreuz{(wA)} \quad \kreuz{(fA)}
\teil $x = -4$ in $x^2 + 16 = 0$ \hfill \kreuz{(wA)} \quad \kreuz{(fA)}
\teil $x = -3$ in $(x - 3)^2 = 36$ \hfill \kreuz{(wA)} \quad \kreuz{(fA)}
\teil $x = 3$ in $(x + 5)^2 = 4$ \hfill \kreuz{(wA)} \quad \kreuz{(fA)}
\teil $x = -2$ in $2x^2 + 1 = 9$ \hfill \kreuz{(wA)} \quad \kreuz{(fA)}
\teil $x = -1$ in $(x - 1)^2 = -4$ \hfill \kreuz{(wA)} \quad \kreuz{(fA)}
\end{teile}
\end{aufgabe}

\begin{aufgabe}{Ich kann am Graphen ablesen, wie viele Lösungen eine Gleichung hat.}
Die Grafik zeigt die Parabel zu $y = (x - 1)^2 - 2$ und drei waagerechte Geraden.

\begin{ksys}[xmin=-3,xmax=5,ymin=-4,ymax=6,ablesen]
\parabel{1}{1}{-2}{}
\gerade[4.4]{0}{2}{y=2}
\gerade[4]{0}{-2}{y=-2}
\gerade[4]{0}{-3}{y=-3}
\end{ksys}

Gib an, wie viele Lösungen die Gleichung hat.
\begin{teile}
\teil $(x - 1)^2 - 2 = 2$ \hfill Anzahl: \leerfeld
\teil $(x - 1)^2 - 2 = -2$ \hfill Anzahl: \leerfeld
\teil $(x - 1)^2 - 2 = -3$ \hfill Anzahl: \leerfeld
\teil $(x - 1)^2 - 2 = 0$ \hfill Anzahl: \leerfeld
\anweisung{Lies die Lösungen am Graphen ab.}
\teil $(x - 1)^2 - 2 = 2$ \hfill \feld{x_1} \quad \feld{x_2}
\end{teile}
\end{aufgabe}

\begin{aufgabe}{Ich kann eine Gleichung angeben, die keine Lösung hat.}
Gib eine passende Gleichung an.
\begin{teile}
\teil Form $x^2 = $ Zahl, ohne Lösung \hfill \leerfeld
\teil mit quadrierter Klammer, ohne Lösung \hfill \leerfeld
\teil mit $x^2$ und einer Zahl links, ohne Lösung \hfill \leerfeld
\end{teile}
\end{aufgabe}

\begin{aufgabe}{Ich kann Aussagen über die Lösungen einer Gleichung beurteilen.}
Kreuze an, ob die Aussage wahr oder falsch ist.
\begin{teile}
\teil Die Gleichung $(x - 5)^2 = 0$ hat genau eine Lösung. \hfill \kreuz{wahr} \quad \kreuz{falsch}
\teil $2$ und $10$ sind die Lösungen der Gleichung $(x + 6)^2 = 16$. \hfill \kreuz{wahr} \quad \kreuz{falsch}
\anweisung{Ändere die Gleichung $(x + 6)^2 = 16$ so, dass sie keine Lösung hat.}
\teil Neue rechte Seite: \leerfeld \quad (P10 2021 FOR)
\end{teile}
\end{aufgabe}

\begin{aufgabe}{Ich finde den Fehler beim Wurzelziehen.}
Finde den Fehler, benenne ihn und korrigiere die Rechnung.
\begin{teile}
\teil Ben löst $2x^2 = 162$:
\rechnung{2x^2 &= 162 &&\mid :2 \\ x^2 &= 81 \\ x &= 9}
\teil Lina löst $(x + 2)^2 = 49$:
\rechnung{x + 2 &= 7 &&\text{oder}\quad x + 2 = -7 \\ x_1 &= 9 &&\text{oder}\quad x_2 = -5}
\teil Tom löst $x^2 + 21 = 12$:
\rechnung{x^2 + 21 &= 12 &&\mid -21 \\ x^2 &= -9 \\ x &= -3}
\end{teile}
\end{aufgabe}

\begin{aufgabe}{Ich kann begründen, warum eine Gleichung zwei, eine oder keine Lösung hat.}
Begründe ohne zu rechnen.
\begin{teile}
\teil Warum hat $x^2 = 20$ zwei Lösungen?
\teil Warum hat $x^2 = -20$ keine Lösung?
\teil Tim sagt: „$(x - 4)^2 = 0$ hat keine Lösung, weil man aus $0$ keine Wurzel ziehen kann.“ Stimmt das?
\end{teile}
\end{aufgabe}

\begin{aufgabe}{Ich kann eine Länge durch Wurzelziehen berechnen.}
Stelle eine Gleichung auf, löse sie und antworte mit Einheit.
\begin{teile}
\teil Eine quadratische Terrasse hat einen Flächeninhalt von $64\,\text{m}^2$. Wie lang ist eine Seite der Terrasse? (P10 2020 FOR) \hfill Seite: \leerfeld[m]
\teil Eine zylinderförmige Regentonne ist $90\,\text{cm}$ hoch und soll $300$ Liter fassen. Für das Volumen gilt $V = \pi\cdot r^2\cdot h$ mit dem Radius $r$ und der Höhe $h$. Wie groß muss der Radius sein? Runde auf eine Stelle nach dem Komma. (P10 2022 FOR) \hfill Radius: \leerfeld[cm]
\end{teile}
\end{aufgabe}
